% !TeX root = ../papers/nc-p3.tex
\section{二次曲線の平面の方程式が与える連続純関係}
\label{sec:conic-corridor-excess}
\begin{proposition}\label{prop:conic-corridor-excess}
曲面射影のない標準正則窓で、順中心 $p$ からの第一直線を $L$、第二曲線を平面二次曲線 $C$、
その平面の方程式を $h_C$ とする。$R_i$ の $L$ への制限は
$\ker p\otimes H^0(L,\mathcal O(1))$ への同型である。
$p$ を通る曲線枝において、順評価の三次小行列式の共通位数が点イデアルの位数を超えるなら、
ある非零 $s\in L^\perp$ に対して $s\otimes h_C\in R_{i+1}$ である。
列の第三次数も2なら、その第三曲線 $C'$ に対して $R_{i+2}$ は第二因子 $h_{C'}$ の純関係を含む。
後者には位数の仮定は不要である。逆方向にも成立する。
\end{proposition}
\begin{proof}
第一・第二点を独立に動かした縮約は $[s,t]A(t)$ と書ける。ここで $A(t)$ は
$L$ 上の唯一の順中心 $q$ の鉛筆で、三行四列の一次行列である
（定理\ref{thm:rank-ten-conic-after-center}の標準二列による同定）。
定理\ref{thm:finite-pure-relation-locus}と命題\ref{prop:finite-pure-pencil-opposite}より
一般階数は3である。中央支持3なら一般右核は固定で、第三点が動かない矛盾となるので
中央支持は4である。一般の $t$ ごとに $A(t)$ の像は $\ker p$ 全体で、
二つの異なる一般 $t$ で $[-,t]$ は独立である。従って制限の像は六次元であり同型を得る。
また第三曲線は $A(t)$ の $L_e$ ブロックの核曲線なので、その次数は $e$ である。
$e=2$ なら残る一行一列が与える一関係はその二項一次式と、残る第四文字との積である。
核曲線は最初の三座標の平面を張るから、第四文字がちょうど $h_{C'}$ である。

命題\ref{prop:excess-contact-two-line-pure}の分割した四行を使う。
制限が同型なので、枝の先頭接方向 $\xi\in(\ker p)^*$ と $L$ の二座標の対が
独立な除算後の二行を与える。過剰な位数があれば、これらは $p$ の残る二行と同じ二空間を張る。
二行を相補切断により $R_{i+2}$ の元として $C$ に制限すると、
符号と基底を除いて $t_0\xi A(t),t_1\xi A(t)$ である。
この二行の末尾の文字の支持は、一次行 $\xi A(t)$ の二係数が張る高々二次元空間に入る。
従って末尾から $w(p)$ を縮約した三空間 $P\subset L^\perp\otimes V_{i+2}$ を
$C$ に制限した像も高々二次元である。核には非零の元があり、
$C$ の線型方程式は $h_C$ の倍だけなので、これが求める純関係である。
\end{proof}

\begin{corollary}\label{cor:conic-corridor-consecutive}
$(1,2,1)$ の列の始点・終点の両方に、それぞれ順・逆の過剰な位数を持つ枝があることはない。
$(1,2,2,1)$ の列では終点に逆の過剰な位数を持つ枝があることはなく、始点についても同様である。
\end{corollary}
\begin{proof}
前者では中央二次曲線の方程式を中間文字として二つの純関係が連続する。
後者では第三曲線の平面の方程式について、順向きの命題の後半と逆向きの前半で同じ矛盾を得る。
連続純関係の禁止は標準正則性の必要条件である。
\end{proof}
